\documentclass[12pt,a4paper]{article}
\usepackage[UTF8,heading=true]{ctex}
\usepackage{geometry}
\usepackage{amsmath,amssymb}
\usepackage{enumitem,hyperref,fancyhdr,setspace}

\geometry{left=2.5cm,right=2.5cm,top=2.5cm,bottom=2.5cm}
\setlength{\headheight}{15pt}
\setstretch{1.5}
\setlength{\emergencystretch}{3em}
\hfuzz=20pt
\sloppy
\hypersetup{hidelinks}

\pagestyle{fancy}
\fancyhf{}
\fancyhead[C]{发明专利权利要求书}
\fancyfoot[C]{\thepage}
\renewcommand{\headrulewidth}{0pt}

\begin{document}

\begin{center}
    {\Large\heiti 中国移动专利申请} \\[0.5em]
    {\Huge\heiti 权利要求书} \\[0.8em]
\end{center}

\noindent\textbf{发明名称：}基于硅基光电集成的微型化 QKD 芯片封装与校准技术

\vspace{0.8em}

\begin{enumerate}[
    label=\arabic*.,
    leftmargin=0pt,
    itemindent=2em,
    labelsep=0.5em,
    align=left,
    itemsep=0.6em,
    topsep=0pt
]

\item 一种面向量子密钥分发 QKD 芯片的硅基光电集成封装结构，其特征在于，包括：
硅光子芯片，所述硅光子芯片上形成承载量子业务的主光路、与所述主光路共模布置的参考波导、用于表征温度漂移的片上传感单元、用于表征封装应力或应变的传感单元以及用于补偿相位、偏置或损耗的执行单元；
封装基座，所述硅光子芯片固定于低热膨胀基座或低应力基座上；
应力释放结构，设置于所述硅光子芯片周围或封装互连区域，用于减小外壳形变、装配残余应力或温度循环对核心光路区的影响；
控制电路，配置为基于所述参考波导和所述传感单元的观测结果生成对所述执行单元的补偿控制量；
其中，所述参考波导与所述主光路在热环境和几何路径上至少部分共模，以提高对温度漂移和应力漂移的可观测性，并且所述控制电路被配置为在量子业务时隙之外或业务允许的校准时隙内执行校准。

\item 如权利要求1所述的封装结构，其特征在于，所述参考波导与所述主光路并排、镜像或环绕布置，并在同一区域内与热源器件保持预设距离，以使共模温漂被共同感知而差分漂移被单独估计。

\item 如权利要求1所述的封装结构，其特征在于，所述片上传感单元包括以下一项或多项：功率抽头与监测光电二极管、参考马赫-曾德干涉仪、参考微环谐振器、温度传感器、应变计以及应力释放梁的电学读出结构。

\item 如权利要求1所述的封装结构，其特征在于，所述执行单元包括以下一项或多项：加热器阵列、相位调制器偏置调节单元、IQ 偏置调节单元、可调衰减器以及光开关。

\item 如权利要求1所述的封装结构，其特征在于，所述应力释放结构包括应力释放槽、柔性连接梁、对称焊盘、对称粘接区或低模量缓冲层中的至少一种，以降低封装材料蠕变、跌落冲击或弯折引起的偏移。

\item 如权利要求1所述的封装结构，其特征在于，所述封装结构进一步包括光学输入输出耦合单元，所述光学输入输出耦合单元采用边耦合、光栅耦合或聚合物波导过渡中的至少一种，并在所述光学输入输出耦合单元周围设置应变隔离结构以提高耦合稳定性。

\item 如权利要求1所述的封装结构，其特征在于，所述控制电路与所述硅光子芯片近距共封装为模组，并通过短互连连接至所述传感单元和所述执行单元，以降低反馈闭环延迟与片外连线引入的噪声。

\item 一种针对权利要求1至7任一项所述封装结构的 QKD 芯片校准控制方法，其特征在于，包括：
将量子业务组织为包含量子业务时隙和静默校准时隙的帧结构或超级帧结构；
在静默校准时隙中开启参考光路，以低占空比、低功率向所述参考波导和选定主光路段注入校准信号；
采集所述参考波导、所述片上传感单元和所述应力传感单元的观测结果；
由控制电路根据所述观测结果估计温度漂移分量与应力漂移分量，并生成对所述执行单元的补偿控制量；
在后续量子业务时隙中使主光路按照补偿后的工作点运行，以维持量子态制备或检测稳定性。

\item 如权利要求8所述的方法，其特征在于，所述静默校准时隙被插入至量子业务帧内部，而非持续注入高功率参考光，从而降低参考信号对单光子探测背景噪声的污染。

\item 如权利要求8所述的方法，其特征在于，所述参考光路与量子业务光路之间的隔离至少通过以下一项或多项实现：时分隔离、波分隔离、片上滤波、封装内滤波、探测窗口门控以及参考信号调制标记识别。

\item 如权利要求8所述的方法，其特征在于，所述估计温度漂移分量与应力漂移分量包括采用卡尔曼滤波、查找表补偿、模型预测控制、在线辨识或数据驱动回归中的至少一种，对状态变量
$\mathbf{x}_k$ 进行更新，并分别输出对应不同执行单元的补偿量。

\item 如权利要求8所述的方法，其特征在于，所述控制方法具有多时间尺度闭环，其中针对毫秒级快速升温、振动冲击或充电事件执行快速补偿，针对秒级温漂或应力蠕变执行深度重锁定或低频扫描。

\item 如权利要求8所述的方法，其特征在于，当检测到参考波导与主光路之间的差分相位偏移超过阈值、微环谐振偏移连续超过阈值持续预设采样周期数、或应力传感输出超过阈值时，触发快速补偿或深度重锁定。

\item 如权利要求8所述的方法，其特征在于，当终端侧事件表明设备处于充电、通话、射频高功率发射、视频拍摄、跌落恢复或运动状态变化时，临时提高静默校准时隙密度或调整控制律参数。

\item 如权利要求8所述的方法，其特征在于，当扰动过大以致无法在预设时隙内恢复目标工作点时，进入降级保护状态，降低密钥业务速率、切换至更稳健工作模式或延后深度重锁定，以保持至少最低可用密钥率。

\item 一种微型化 QKD 模组，其特征在于，包括如权利要求1至7任一项所述的封装结构以及被配置为执行如权利要求8至15任一项所述方法的控制电路，其中所述微型化 QKD 模组被封装为可嵌入手机、智能手表、便携终端、小型网关或加密钥匙的器件。

\item 如权利要求16所述的微型化 QKD 模组，其特征在于，所述微型化 QKD 模组通过 I\textsuperscript{2}C、SPI 或高速串行接口与主控芯片连接，并向主控芯片输出温度状态、应力状态、锁定状态和告警信息。

\item 一种终端设备，其特征在于，包括壳体、主控芯片、电源管理单元以及如权利要求16或17所述的微型化 QKD 模组，所述主控芯片向所述微型化 QKD 模组提供终端事件信息，所述微型化 QKD 模组根据所述终端事件信息调整静默校准时隙密度与控制策略。

\item 一种电子设备，包括存储器、处理器及存储在所述存储器中并可在所述处理器上运行的计算机程序，其特征在于，所述处理器执行所述计算机程序时实现如权利要求8至15任一项所述的方法。

\item 一种计算机可读存储介质，其上存储有计算机程序，其特征在于，所述计算机程序被处理器执行时实现如权利要求8至15任一项所述的方法。

\end{enumerate}

\end{document}
